\subsection{The exact once-punctured elliptic comparison in intrinsic class three}

\begin{statement}{E.50}{Proved}{The multiplication-by-three cover and its integral ambient-metabelian kernel}
Let $E$ be the cubic elliptic curve used in E.37, with
origin $O$, and put $U=E\setminus E[3]$ and
$E^\circ=E\setminus\{O\}$. Choose the actual rational
tangent $v$ at $O$ used for $U$, and use $3v$ as
the tangent on $E^\circ$. The morphism
\[
 [3]:(U,v)\longrightarrow(E^\circ,3v)
 \tag{E50.1}
\]
induces an inclusion of geometric pro-three groups
\[
 N=\pi_1(U_{\overline{\mathbb Q}},v)^{(3)}
 \lhook\joinrel\longrightarrow
 F=\pi_1(E^\circ_{\overline{\mathbb Q}},3v)^{(3)},
 \qquad N=\ker(F\longrightarrow F^{\rm ab}/3),
 \tag{E50.2}
\]
equivariant for the resulting based $G_{\mathbb Q}$
actions. Here $F$ is free pro-three of rank two and
$N$ is free pro-three of rank ten. All derived and
lower central subgroups below are closed.

Write $C=N/\gamma_4N$. The actual map to the image
in the class-three quotient of the ambient
metabelian group is
\[
 C\longrightarrow
 \overline C=(N/F'')/\gamma_4(N/F'')
       =N/(F''\gamma_4N),
 \qquad
 K_0=\ker(C\longrightarrow\overline C)
       =F''\gamma_4N/\gamma_4N.
 \tag{E50.3}
\]
Let $V=N^{\rm ab}$ and let $p:V\to T_3E$ be
the map induced by filling the nine punctures.
The map $V\to F^{\rm ab}=T_3E$ in (E50.2) is
exactly $3p$. Thus
\[
 T=\ker(p)=\ker(V\longrightarrow F^{\rm ab})
 \tag{E50.4}
\]
is the primitive rank-eight cusp lattice of E.37.
Choose the free basis described in the proof,
and write $V=L\oplus T$, where $L$ is spanned by
its two handle generators. This splitting is a
choice of integral coordinates. Then
\[
 \begin{array}{c|ccc}
  &\operatorname{gr}_1&\operatorname{gr}_2&\operatorname{gr}_3\\ \hline
 C&V&\Lambda^2V&\operatorname{Lie}_3(V)\\
 K_0&0&\Lambda^2T&[T,\operatorname{Lie}_2(V)]\\
 \overline C&V&\Lambda^2L\oplus(L\otimes T)&
      \operatorname{Lie}_3(L)\oplus
                      (\operatorname{Sym}^2L\otimes T)
 \end{array}
 \tag{E50.5}
\]
is an exact description of the induced lower
central filtrations over $\mathbb Z_3$. The last
row has ranks $10,17,26$; the kernel row has
ranks $0,28,304$. In particular $K_0$ is an
abelian pro-three group isomorphic to
$\mathbb Z_3^{332}$, and all the displayed
graded quotients are torsion-free. The maps and
kernel in (E50.3) are $G_{\mathbb Q}$-equivariant;
the direct sums involving the chosen $L$ record
the stated coordinate splitting.

There is also a direct ambient route retaining
the whole intrinsic class-three group:
\[
 F'''\subset\gamma_4N,\qquad
 F/F'''\longrightarrow F/\gamma_4N,
 \qquad
 \ker(F/\gamma_4N\longrightarrow F/N)=C.
 \tag{E50.6}
\]
Thus an actual based action on the indicated
ambient derived-length-three quotient induces
the actual action on $C$ through these maps.

The comparison has a concrete automorphism
kernel. There are based automorphisms of $F$
preserving its peripheral conjugacy class which
act identically on $F/F''$ and nontrivially on
$C$. One is conjugation by
$w=[t_1,t_2]\in F''$, in the basis below:
\[
 \operatorname{Ad}_w(e_1)e_1^{-1}
       =[[t_1,t_2],e_1]\ne1\quad\text{in }C.
 \tag{E50.7}
\]
This is an assertion about these specified
based automorphism maps.

Two further transfers retain the exact
geometric and arithmetic origins. The full
peripheral kernel of $N\to T_3E$ is
$I=F'$, and the quotient in E.36 fits into
the actual equivariant sequence
\[
 1\longrightarrow\mathbb Z_3^{42}
 \longrightarrow\overline C
 \longrightarrow N^{\rm cc}=N/[N,I]
 \longrightarrow1.
 \tag{E50.8}
\]
The kernel has induced graded ranks
$0,16,26$ and is abelian; the displayed
isomorphism with $\mathbb Z_3^{42}$ is
an assertion about its underlying group.

Let $\rho_E:G_{\mathbb Q}\to
\operatorname{Aut}(T_3E)$ be the actual
elliptic Tate representation. The based
action on $\overline C$ induces a
continuous homomorphism
\[
 \rho_E(G_{\mathbb Q})
       \longrightarrow\operatorname{Out}(\overline C)
 \tag{E50.9}
\]
whose composition with $\rho_E$ is the
genuine outer action on this exact
quotient. For $g\in\ker\rho_E$ the proof
gives its actual inner conjugator in
$F'/F''\subset N/F''$.

The exact integer Hall certificates are saved
in \texttt{a\_elliptic\_kernel.py}, its output
and JSON file. They verify the cover relation,
1,000 tensor identities, all 304 integral kernel
coordinates, and 280 normality identities.

\emph{Manuscript: \S51.16, folios 733v--736v,
scans 172--178, (371)--(379), Part E(c),(d).
This comparison supplies the precise quotient
needed when applying once-punctured elliptic
formulae to the intrinsic cubic group.}
\end{statement}
\begin{proof}
\emph{The actual cover and the based action.}
Multiplication by three is an etale isogeny in
characteristic zero. Its inverse image of $O$
is $E[3]$, so (E50.1) is a finite etale cover
of degree nine. Its differential at $O$ sends
$v$ to $3v$; these are actual tangential
basepoints, defined over $\mathbb Q$. The
geometric covering subgroup is the kernel of
the homology reduction modulo three. Its
induced pro-three topology is the full one:
for a normal subgroup of three-power index in
the covering group, intersect its nine
conjugates. The resulting quotient of the
covering group is a three-group and the
quotient of the ambient group is an extension
of it by $(\mathbb Z/3)^2$. Taking these
quotients proves the asserted injection of
pro-three completions. The map of tangential
arithmetic fundamental groups gives the
equivariance, including the chosen sections
of $G_{\mathbb Q}$.

Cut a complex torus along a symplectic pair of
loops $x_1,x_2$ based just before the origin,
and let $z=[x_1,x_2]^{-1}$ be its positive
peripheral loop. In the nine-square cover
choose the peripheral paths in the following
order of pairs:
\[
 (0,0),(0,1),(1,1),(1,0),(0,2),
 (1,2),(2,2),(2,1),(2,0).
 \tag{E50.10}
\]
For a pair $(a,b)$ take
$x_{ab}=x_1^ax_2^b$ for $a\ge b$ and
$x_{ab}=x_1^bx_2^bx_1^{a-b}$ for $a<b$,
and put $z_{ab}=x_{ab}zx_{ab}^{-1}$.
The boundary of the nine-square polygon is
the exact relation
\[
 [x_1^3,x_2^3]z_0z_1\cdots z_8=1.
 \tag{E50.11}
\]
These paths and the same relation are given
in Nakamura, \S2.4.1--2.4.2, p.205
\cite{A_Nakamura1995}. The program also
checks (E50.11) directly by free reduction
in $x_1,x_2$. The punctured surface has
genus one and nine punctures, hence its
group is free of rank ten, with basis
\[
 e_1=x_1^3,\quad e_2=x_2^3,
 \quad t_i=z_i\ (1\le i\le8),
 \tag{E50.12}
\]
after eliminating $z_0$ with (E50.11).
The same is true of the pro-three group.
Transporting the complex basepoint to the
tangent gives this geometric basis of the
already specified based group; it does not
alter the arithmetic morphism (E50.1).

Filling the punctures sends the two handle
classes to a basis of $T_3E$ and each $t_i$
to zero. The proper map in the commutative
filling square is $[3]:E\to E$, whose Tate
module map is multiplication by three.
This proves (E50.4) and the precise factor
$3p$. The image of $N$ in $F^{\rm ab}$ is
$3F^{\rm ab}$ and its kernel is $F'$.
Consequently, inside $N$,
\[
 F'=\ker(N\longrightarrow3F^{\rm ab})
    =\langle\!\langle t_1,\ldots,t_8,
                   [e_1,e_2]\rangle\!\rangle_N.
 \tag{E50.13}
\]
Indeed, quotienting by the displayed normal
generators gives precisely the abelian free
pro-three group on $e_1,e_2$, and the latter
maps isomorphically onto $3F^{\rm ab}$.

\emph{Passage to the actual class-three group.}
Integral Hall collection gives unique
$\mathbb Z_3$ coordinates for $C$ on its
ten weight-one generators, its 45 basic
commutators of weight two, and its 330 Hall
commutators of weight three. These are
group coordinates, not rational
Campbell--Hausdorff coordinates. For clarity,
order the ten generators by the indices
$0,1,2,\ldots,9$, with the first two the
handles. The weight-three basis is
\[
 [[v_i,v_j],v_k],\qquad
        i>j,\quad k\ge j.
 \tag{E50.14}
\]
Unique collection for integral free
nilpotent groups, followed by the inverse
limit of their three-power coordinate
quotients, gives these coordinates in the
free pro-three quotient. Equivalently,
the successive central extensions have
the integral Hall lattices constructed
in E.41. In particular this argument
uses no division by three.

Let $J$ be the image of $F'$ in $C$, let
$H$ be the closed subgroup generated by
the $t_i$, and put $A=\gamma_2C$. From
(E50.13) and the abelianization of $C$,
\[
 J=H A,\qquad [A,A]=1,\qquad
                  [H,A]\subset\gamma_3C.
 \tag{E50.15}
\]
In the last inclusion the target is
central. Let $B$ be the closed subgroup
generated by the 28 actual basic group
commutators $[t_i,t_j]$, $i>j$, and let
$L_3$ be the closed subgroup of
$\gamma_3C$ generated by
$[t_i,[v_a,v_b]]$. Commutator collection
in (E50.15) gives
\[
 K_0=J'=B L_3.
 \tag{E50.16}
\]
Here $[H,H]$ is generated by $B$ and
triple commutators on the $t_i$, while
$[H,A]=L_3$. Those triple commutators
already belong to $L_3$. Conversely,
both $B$ and $L_3$ consist of commutators
of elements of $J$, since $A\subset J$.
This proves equality of actual closed
subgroups, before considering their
associated graded modules.

All generators of $B$ lie in $\gamma_2C$,
so they commute with each other, and
$L_3$ is central. Their degree-two Hall
coordinates are 28 distinct basis
vectors. The continuous map
$\mathbb Z_3^{28}\to B$ given by their
ordered powers is therefore injective,
is surjective by the definition of $B$,
and satisfies $B\cap\gamma_3C=1$.
Thus
\[
 K_0=B\times L_3,
 \qquad B\simeq\mathbb Z_3^{28},
 \qquad K_0\cap\gamma_3C=L_3.
 \tag{E50.17}
\]
This supplies the required passage from
graded information to the whole kernel
at the prime three.

\emph{The saturated degree-three lattice.}
In the central group $\gamma_3C$, a
commutator $[t,[v,w]]$ is exactly the
corresponding integral Lie bracket.
Hence $L_3=[T,\operatorname{Lie}_2(V)]$.
The degree-two quotient is obtained by
deleting the basis $\Lambda^2T$.
For degree three, use the decomposition
$V=L\oplus T$. Jacobi shows that
$[V,[T,T]]\subset[T,\operatorname{Lie}_2(V)]$.
Every Lie monomial with at least two
$T$ letters therefore vanishes in the
quotient. For a fixed $t\in T$, the
remaining relation with one $T$ letter
is $[t,[e_1,e_2]]=0$. Consequently
\[
 [e_i,[e_j,t]]=[e_j,[e_i,t]],
 \tag{E50.18}
\]
and the surviving one-$T$ coordinates
are the three symmetric pairs of the
two handles. In the Hall basis
(E50.14), retain exactly the two pure
handle triples and
\[
 [[t,e_1],e_1],\quad
 [[t,e_1],e_2],\quad
 [[t,e_2],e_2]
 \quad(t=t_1,\ldots,t_8).
 \tag{E50.19}
\]
Projection to these 26 coordinates
annihilates all generators of $L_3$.
Conversely, the 304 discarded Hall
coordinates are integral combinations
of those generators. One direct
certificate uses the identity
\[
 [[v_i,v_j],v_k]
   =[[v_i,v_k],v_j]-[[v_j,v_k],v_i]
                  \quad(k<j<i).
 \tag{E50.20}
\]
The included program forms the 360
columns $[t_i,[v_a,v_b]]$ and saves an
integer right inverse onto every
discarded Hall coordinate. All its
nonzero coefficients are $1$ or $-1$;
it verifies the right-inverse matrix
identity over $\mathbb Z$. This is also
an explicit certificate of the
integral generation argument just
given. The decomposition by the
number of $T$ letters is
\[
 \begin{array}{c|rrrr}
 \text{number of $T$ letters}&0&1&2&3\\ \hline
 \operatorname{Lie}_3(V)&2&32&128&168\\
 L_3&0&8&128&168\\
 \operatorname{Lie}_3(V)/L_3&2&24&0&0.
 \end{array}
 \tag{E50.21}
\]
Thus $L_3$ is saturated of rank 304.
Equations (E50.17) and (E50.21) prove
the actual abelian kernel rank 332.
They also give (E50.5) with free
integral quotients. Normality follows
either from (E50.16) or from the
integral Jacobi inclusions; the latter
are checked for every handle and cusp
generator in the program.

\emph{The further ambient map and the
automorphism comparison.}
Since $F'\subset N$, one has
$F''\subset\gamma_2N$ and then
$F'''=[F'',F'']\subset\gamma_4N$.
All these subgroups are preserved by
every automorphism of $F$, since
$N=\ker(F\to F^{\rm ab}/3)$ is
characteristic. This proves (E50.6)
and its compatibility with based
actions. For (E50.7), $w\in F''$
acts trivially on $F/F''$, whereas its
commutator with $e_1$ is a nonzero
Hall combination in $L_3$. Inner
automorphisms preserve the peripheral
conjugacy class. This proves the
specified based-automorphism claim.

For outer actions, changing an
ambient representative by
$\operatorname{Ad}_f$ changes its
restriction by the action of $f$ on
$N$. Elements of $N$ give inner
automorphisms of $N$; the remaining
ambiguity is the deck action of
$F/N=(\mathbb Z/3)^2$. The tangent
map in (E50.1) fixes the based
representatives used here. This
describes the precise transfer
without choosing an unspecified
outer lift.

\emph{The transfer to the cuspidally
central quotient.}
The filling map has kernel $I$ and
its composite with multiplication
by three has kernel $F'$ by
(E50.13). Multiplication by three
is injective on $T_3E$, so $I=F'$.
Consequently
\[
 F''=[I,I]\subset[N,I],\qquad
 \gamma_3N=[N,N']\subset[N,I],
 \tag{E50.22}
\]
and (E50.8) is the resulting
quotient map. In the class-three
Hall coordinates, $[N,I]$ is
generated by the weight-two
commutators with at least one
$T$ letter and all of
$\gamma_3C$. Modulo $K_0$ its
weight-two lattice is $L\otimes T$
of rank sixteen and its central
weight-three lattice is the
rank-twenty-six quotient in
(E50.5). These are subsets of
the remaining Hall coordinates.
Their actual group commutators
vanish, because they have weight
at least four. As in (E50.17),
independent weight-two coordinates
split off the central part. This
proves the abelian kernel and its
rank forty-two. The target and
its arithmetic origin are exactly
the group $N/[N,I]$ of E.36.

\emph{The actual outer factorization.}
Use $z=[x_1,x_2]^{-1}$ in $F$ and
let
\[
 \mathcal A=\mathbb Z_3[[F^{\rm ab}]]
      =\mathbb Z_3[[T_1,T_2]],\qquad
 X_i=1+T_i.
 \tag{E50.23}
\]
The two-generator instance of the
completed Fox resolution proved
in E.41 identifies $F'/F''$ with
the free $\mathcal A$-module on
$\bar z$: the syzygy kernel of
$(a,b)\mapsto aT_1+bT_2$ is
generated by $(T_2,-T_1)$, the
Fox image of $\bar z$.

Fix $g\in\ker\rho_E$ and let $f_g$
be its actual tangential action
reduced on $F/F''$. The Tate
determinant is $\chi_g$, so
$\chi_g=1$. Tangential inertia
therefore gives $f_g(\bar z)=\bar z$
exactly, while $f_g$ acts trivially
on the abelianization. Write
\[
 f_g(\bar x_i)\bar x_i^{-1}
      =G_i\cdot\bar z\quad(i=1,2).
 \tag{E50.24}
\]
Computing the commutator of these
two images in the metabelian
group gives
$f_g(\bar z)=(1+T_2G_1-T_1G_2)
\cdot\bar z$. Thus
$T_2G_1=T_1G_2$. Reduction modulo
$T_1$ shows $G_1\in(T_1)$,
because $T_2$ is a non-zero-divisor
in $\mathcal A/(T_1)$. Write
$G_1=T_1H_g$ and cancel $T_1$
to obtain $G_2=T_2H_g$.
This determines $H_g$ uniquely
in the actual integral ring.
Conjugation by
\[
 w_g=-H_g\cdot\bar z\in F'/F''
                           \subset N/F''
 \tag{E50.25}
\]
changes $\bar x_i$ by
$(1-X_i)w_g=T_iH_g\cdot\bar z$.
It therefore equals $f_g$ on the
two generators and hence on the
whole group. This is the direct
calculation in Nakamura's
Proposition4.7--\S4.8
\cite{A_Nakamura1995}, applied
to the actual tangential lift.
Its restriction and quotient
on $\overline C$ is inner by
the image of precisely $w_g$.
The genuine outer action thus
kills $\ker\rho_E$. Since
$G_{\mathbb Q}$ is compact
and its Tate image is
Hausdorff, descent gives the
continuous homomorphism
(E50.9). This proves the
factorization in its stated
outer automorphism group.

Finally, in Nakamura's \S6.3
$I_m=\ker(\mathbb Z_3[[F]]\to
\mathbb Z_3[[F/\gamma_mF]])$ differs
from the augmentation power $I^m$.
The Magnus formula in Theorem6.6
and Corollary6.7 is modulo $I_2$;
its full noncommutative matrix is
defined in \S6.5. The maps
(E50.3) and (E50.6) identify the
respective intrinsic targets for
that formula and for an actual
ambient derived-length-three
action \cite{A_Nakamura1995}.
\end{proof}
